\chapter{The Transforms of Applied Mathematics}

\section{Finite Fourier Transform (Fourier Series) of Periodic Functions}

\begin{definition}{Periodic Functions and the Complex Fourier Series}{fourier_series_def}
A function $F(t)$ is \textbf{periodic} with period $L > 0$ if
\begin{equation}
F(t + L) = F(t) \quad \forall t \in \mathbb{R}.
\end{equation}
The discrete fundamental frequencies are defined by:
\begin{equation}
\omega_n := \frac{2\pi n}{L}, \quad n \in \mathbb{Z}.
\end{equation}
If $F(t)$ is integrable over $\left[-\frac{L}{2}, \frac{L}{2}\right]$, its \textbf{complex Fourier series} is:
\begin{equation}
F(t) = \sum_{n=-\infty}^\infty c_n e^{i \omega_n t} = \sum_{n=-\infty}^\infty c_n e^{i \frac{2\pi n}{L} t},
\end{equation}
where the \textbf{Fourier coefficients} $c_n$ are given by:
\begin{equation}
c_n = \frac{1}{L} \int_{-L/2}^{L/2} F(t) e^{-i \omega_n t}\, dt, \quad n \in \mathbb{Z}.
\end{equation}
\end{definition}

\begin{proposition}{Rescaled Fourier Series ($L = 2\pi$)}{rescaled_fourier}
When normalized to period $L = 2\pi$, the angular frequencies are integers $\omega_n = n$, and:
\begin{equation}
F(t) = \sum_{n=-\infty}^\infty c_n e^{int}, \qquad c_n = \frac{1}{2\pi} \int_{-\pi}^\pi F(t) e^{-int}\, dt.
\end{equation}
\end{proposition}

\begin{remark}[Connection between Fourier Series and Laurent Series]
Setting $z = e^{it}$, the unit circle $|z| = 1$ is parametrized by $t \in [-\pi, \pi]$. The Fourier series of $F(t)$ corresponds precisely to the Laurent series of an analytic function $f(z)$ on the unit circle:
\begin{equation}
f(z) = \sum_{n=-\infty}^\infty c_n z^n \implies F(t) = f(e^{it}) = \sum_{n=-\infty}^\infty c_n e^{int}.
\end{equation}
The Fourier coefficients $c_n$ match the Laurent contour integral coefficients:
\begin{equation}
c_n = \frac{1}{2\pi i} \oint_{|z|=1} \frac{f(z)}{z^{n+1}}\, dz = \frac{1}{2\pi} \int_{-\pi}^\pi F(t) e^{-int}\, dt.
\end{equation}
\end{remark}

\begin{theorem}{Convergence of Fourier Series (Dirichlet Conditions)}{fourier_convergence}
Let $F(t)$ be a piecewise smooth periodic function with period $L$. Then the Fourier series of $F(t)$ converges for all $t$:
\begin{equation}
\sum_{n=-\infty}^\infty c_n e^{i \omega_n t} = \frac{F(t^+) + F(t^-)}{2}.
\end{equation}
At every point $t$ where $F$ is continuous, the series converges directly to $F(t)$.
\end{theorem}

\begin{proposition}{Termwise Differentiation of Fourier Series}{fourier_diff}
If $F(t)$ is continuous and piecewise smooth with period $L$, its derivative $F'(t)$ has the Fourier series obtained by term-by-term differentiation:
\begin{equation}
F'(t) = \sum_{n=-\infty}^\infty (i \omega_n) c_n e^{i \omega_n t}.
\end{equation}
\end{proposition}

\begin{theorem}{Parseval's Identity for Fourier Series}{parsevals_identity}
If $F(t)$ is square-integrable over a period, the total energy is conserved:
\begin{equation}
\frac{1}{L} \int_{-L/2}^{L/2} |F(t)|^2\, dt = \sum_{n=-\infty}^\infty |c_n|^2.
\end{equation}
\end{theorem}


\section{The Fourier Transform}

\begin{definition}{The Fourier Transform and Inversion Formula}{fourier_transform_def}
Let $F(t)$ be a non-periodic, integrable function on $(-\infty, \infty)$.
\begin{itemize}
    \item The \textbf{Fourier transform} of $F(t)$, denoted $\mathcal{F}[F](\omega)$ or $G(\omega)$, is defined by:
    \begin{equation}
    G(\omega) \equiv \mathcal{F}[F](\omega) := \frac{1}{2\pi} \int_{-\infty}^\infty F(t) e^{-i\omega t}\, dt.
    \end{equation}
    \item The \textbf{Fourier inversion formula} reconstructs $F(t)$ from its transform:
    \begin{equation}
    F(t) = \int_{-\infty}^\infty G(\omega) e^{i\omega t}\, d\omega.
    \end{equation}
    At points of jump discontinuity, the integral converges to $\frac{F(t^+) + F(t^-)}{2}$.
\end{itemize}
\end{definition}

\begin{proposition}{Evaluation via Residue Calculus and Jordan's Lemma}{fourier_residues}
Fourier integrals of rational functions $F(t) = \frac{P(t)}{Q(t)}$ can be evaluated using complex contour integration:
\begin{equation}
G(\omega) = \frac{1}{2\pi} \int_{-\infty}^\infty F(t) e^{-i\omega t}\, dt.
\end{equation}
By Jordan's Lemma:
\begin{itemize}
    \item If $\omega < 0$, the factor $e^{-i\omega z} = e^{-i\omega(x+iy)} = e^{-i\omega x} e^{\omega y}$ decays in the \textbf{upper half-plane} ($\Im(z) > 0$). We close the contour with an upper semicircle $C_R^+$:
    \begin{equation}
    G(\omega) = i \sum_{\Im(z_k) > 0} \Res\big(F(z) e^{-i\omega z}; z_k\big) \quad (\omega < 0).
    \end{equation}
    \item If $\omega > 0$, the factor decays in the \textbf{lower half-plane} ($\Im(z) < 0$). We close the contour with a lower semicircle $C_R^-$ traversed clockwise:
    \begin{equation}
    G(\omega) = -i \sum_{\Im(z_k) < 0} \Res\big(F(z) e^{-i\omega z}; z_k\big) \quad (\omega > 0).
    \end{equation}
\end{itemize}
\end{proposition}


\section{The Laplace Transform}

\subsection{Definition and Convergence}

\begin{definition}{The Laplace Transform}{laplace_transform_def}
Let $F(t)$ be defined for $t \ge 0$. The \textbf{Laplace transform} of $F(t)$, denoted $\mathcal{L}[F](s)$, is defined by:
\begin{equation}
\mathcal{L}[F](s) \equiv F(s) := \int_0^\infty F(t) e^{-st}\, dt,
\end{equation}
for all complex numbers $s = \sigma + i\omega$ for which the improper integral converges.
\end{definition}

\begin{proposition}{Condition for Existence (Functions of Exponential Order)}{exponential_order}
If $F(t)$ is piecewise continuous on every finite interval $[0, T]$ and of \textbf{exponential order} $\alpha$, meaning there exist constants $M > 0$ and $\alpha \in \mathbb{R}$ such that
\begin{equation}
|F(t)| \le M e^{\alpha t} \quad \forall t \ge 0,
\end{equation}
then the Laplace transform $\mathcal{L}[F](s)$ exists and is an \textbf{analytic function} in the half-plane:
\begin{equation}
\Re(s) > \alpha.
\end{equation}
\end{proposition}

\subsection{Table of Common Laplace Transforms}

\begin{table}[h!]
\centering
\renewcommand{\arraystretch}{1.5}
\begin{tabular}{|l|l|l|}
\hline
\textbf{Function $F(t)$ ($t \ge 0$)} & \textbf{Laplace Transform $\mathcal{L}[F](s)$} & \textbf{Region of Convergence} \\ \hline
$1$ & $\dfrac{1}{s}$ & $\Re(s) > 0$ \\ \hline
$t^n \quad (n \in \mathbb{N})$ & $\dfrac{n!}{s^{n+1}}$ & $\Re(s) > 0$ \\ \hline
$e^{at} \quad (a \in \mathbb{C})$ & $\dfrac{1}{s - a}$ & $\Re(s) > \Re(a)$ \\ \hline
$\cos(\omega t)$ & $\dfrac{s}{s^2 + \omega^2}$ & $\Re(s) > 0$ \\ \hline
$\sin(\omega t)$ & $\dfrac{\omega}{s^2 + \omega^2}$ & $\Re(s) > 0$ \\ \hline
$\cosh(at)$ & $\dfrac{s}{s^2 - a^2}$ & $\Re(s) > |\Re(a)|$ \\ \hline
$\sinh(at)$ & $\dfrac{a}{s^2 - a^2}$ & $\Re(s) > |\Re(a)|$ \\ \hline
$e^{at} F(t)$ & $\mathcal{L}[F](s - a)$ & $\Re(s) > \alpha + \Re(a)$ \\ \hline
\end{tabular}
\caption{Common Laplace Transforms.}
\end{table}

\subsection{Transform of Derivatives}

\begin{theorem}{Laplace Transform of Derivatives}{laplace_derivatives}
Let $F(t)$ be continuous and of exponential order $\alpha$, with piecewise continuous derivative $F'(t)$. Then:
\begin{equation}
\mathcal{L}[F'](s) = s \mathcal{L}[F](s) - F(0^+).
\end{equation}
For higher-order derivatives:
\begin{align}
\mathcal{L}[F''](s) &= s^2 \mathcal{L}[F](s) - s F(0^+) - F'(0^+), \\
\mathcal{L}[F^{(n)}](s) &= s^n \mathcal{L}[F](s) - \sum_{k=1}^n s^{n-k} F^{(k-1)}(0^+).
\end{align}
\end{theorem}

\subsection{Inverse Laplace Transform via the Bromwich Integral}

\begin{theorem}{Bromwich Inversion Formula}{bromwich_integral}
Let $F(t)$ be piecewise smooth and of exponential order $\alpha$ ($|F(t)| \le M e^{\alpha t}$).
For any real number $a > \alpha$, the inverse Laplace transform is given by the \textbf{Bromwich integral}:
\begin{equation}
\frac{F(t^+) + F(t^-)}{2} = \frac{1}{2\pi i} \text{P.V.} \int_{a - i\infty}^{a + i\infty} \mathcal{L}[F](s) e^{st}\, ds \quad (t > 0).
\end{equation}
\end{theorem}

\begin{theorem}{Evaluation of Inverse Laplace Transforms via Residues}{bromwich_residues}
Let $\mathcal{L}[F](s)$ be analytic in the complex plane except for isolated singularities $s_1, s_2, \dots, s_m$ all lying to the left of the vertical line $\Re(s) = a$.
If $|\mathcal{L}[F](s)| \le \frac{M}{|s|^k}$ ($k > 0$) as $|s| \to \infty$ in the left half-plane $\Re(s) \le a$, then closing the contour to the left with a large semicircle gives:
\begin{equation}
F(t) = \sum_{k=1}^m \Res\big(\mathcal{L}[F](s) e^{st}; s_k\big) \quad \forall t > 0.
\end{equation}
For $t < 0$, closing the contour to the right (which encloses no singularities) yields $F(t) = 0$.
\end{theorem}

\begin{example}{Inverse Laplace Transform via Residues}{inv_laplace_ex}
To find the inverse transform of $\mathcal{L}[F](s) = \frac{1}{s^4 - 1}$:
The function has four simple poles at $s = 1, -1, i, -i$.
The residue of $\frac{e^{st}}{s^4 - 1}$ at a simple root $s_k$ is $\frac{e^{s_k t}}{4 s_k^3} = \frac{s_k e^{s_k t}}{4}$:
\begin{align}
F(t) &= \sum_{k} \Res\left(\frac{e^{st}}{s^4 - 1}; s_k\right) \nonumber \\
     &= \frac{e^t}{4(1)^3} + \frac{e^{-t}}{4(-1)^3} + \frac{e^{it}}{4(i)^3} + \frac{e^{-it}}{4(-i)^3} \nonumber \\
     &= \frac{e^t - e^{-t}}{4} - \frac{e^{it} - e^{-it}}{4i} \nonumber \\
     &= \frac{\sinh t - \sin t}{2} \quad (t \ge 0).
\end{align}
\end{example}

